\documentclass[../../main.tex]{subfiles}

\begin{document}

\section{习题}\label{section:ch12-习题}

以下习题集包括线性和非线性波动方程的问题，以及相关的非线性 Schr\"odinger 方程。除非另有说明，所有给定函数假设光滑。

\begin{exercise}\label{exercise:ch12-能量}
假设 $u$ 具有空间紧支集且求解拟线性波动方程
\begin{align*}
u_{tt}-\sum_{i=1}^n(L_{p_i}(Du))_{x_i}=0 \quad \text{在 } \mathbb{R}^n\times(0,\infty) \text{ 中}.
\end{align*}
确定适当的能量 $E(t)$ 并证明 $\dot{E}\equiv 0$。
\end{exercise}

\begin{exercise}\label{exercise:ch12-KleinGordon}
设 $u$ 求解 Klein--Gordon 方程
\begin{align}
\begin{cases}
u_{tt}-\Delta u+m^2u=0 & \text{在 } \mathbb{R}^n\times(0,\infty) \text{ 中},\\
u=g,\;u_t=h & \text{在 } \mathbb{R}^n\times\{t=0\} \text{ 上}.
\end{cases}
\label{align:ch12-KG}
\end{align}
\begin{enumerate}[(1)]
\item 证明能量
\begin{align*}
E(t)\coloneqq\frac{1}{2}\int_{\mathbb{R}^n}u_t^2+|Du|^2+m^2u^2\,\mathrm{d}x \quad (t\geqslant 0)
\end{align*}
在时间上是常数。
\item 修改 \S4.3.1 中证明波动方程能量渐近均分的论证，证明
\begin{align*}
\lim_{t\to\infty}\int_{\mathbb{R}^n}|Du|^2+m^2u^2\,\mathrm{d}x=E(0).
\end{align*}
\end{enumerate}
\end{exercise}

\begin{exercise}\label{exercise:ch12-KG降维}
假设 $u$ 求解\cref{exercise:ch12-KleinGordon}中 Klein--Gordon 方程的初值问题\eqref{align:ch12-KG}。写 $\bar{x}=(x,x_{n+1})$（$x\in\mathbb{R}^n$），定义
\begin{align*}
\bar{u}(\bar{x},t)\coloneqq u(x,t)\cos(mx_{n+1}).
\end{align*}
\begin{enumerate}[(1)]
\item 证明 $\bar{u}$ 求解波动方程 $\square\bar{u}=0$ 在 $\mathbb{R}^{n+1}\times(0,\infty)$ 中。
\item 推导 $n=1$ 时 Klein--Gordon 方程初值问题解的公式。（这是 \S2.4.1 中引入的降维法的变体。）
\end{enumerate}
\end{exercise}

\begin{exercise}\label{exercise:ch12-阻尼}
假设 $u$ 求解
\begin{align*}
u_{tt}-\Delta u+du_t=0 \quad \text{在 } \mathbb{R}^n\times(0,\infty) \text{ 中},
\end{align*}
对 $d>0$ 是阻尼波动方程。找一个简单指数项，当乘以 $u$ 时给出
\begin{align*}
v_{tt}-\Delta v+cv=0
\end{align*}
的解 $v$，常数 $c<0$。（这与 Klein--Gordon 方程的符号相反。）
\end{exercise}

\begin{exercise}\label{exercise:ch12-相速度}
验证对每个给定的 $y\in\mathbb{R}^n$，$y\neq 0$，函数 $u=e^{\mathrm{i}(x\cdot y-\sigma t)}$ 求解 Klein--Gordon 方程
\begin{align*}
u_{tt}-\Delta u+m^2u=0
\end{align*}
只要 $\sigma=(|y|^2+m^2)^{1/2}$。此平面波解的相速度为 $\sigma/|y|>1$。为什么这不与 \S12.1 中解的传播速度小于或等于 $1$ 的断言矛盾？
\end{exercise}

\begin{exercise}\label{exercise:ch12-三维衰减}
假设
\begin{align*}
\begin{cases}
u_{tt}-\Delta u=0 & \text{在 } \mathbb{R}^3\times(0,\infty) \text{ 中},\\
u=g,\;u_t=h & \text{在 } \mathbb{R}^3\times\{t=0\} \text{ 上},
\end{cases}
\end{align*}
其中 $g,h$ 具有紧支集。证明存在常数 $C$ 使得
\begin{align*}
|u(x,t)|\leqslant C/t \quad (x\in\mathbb{R}^3,\;t>0).
\end{align*}
（提示：使用 \S2.4.1 中解的表示公式。）
\end{exercise}

\begin{exercise}\label{exercise:ch12-二维衰减}
设 $u$ 求解
\begin{align*}
\begin{cases}
u_{tt}-\Delta u=0 & \text{在 } \mathbb{R}^2\times(0,\infty) \text{ 中},\\
u=g,\;u_t=h & \text{在 } \mathbb{R}^2\times\{t=0\} \text{ 上},
\end{cases}
\end{align*}
其中 $g,h$ 具有紧支集。证明
\begin{align*}
|u(x,t)|\leqslant C/t^{1/2} \quad (x\in\mathbb{R}^2,\;t>0)
\end{align*}
对某常数 $C$。
\end{exercise}

\begin{exercise}\label{exercise:ch12-二维爆破}
假设 $u^\varepsilon$ 求解 $n=2$ 空间维数中的线性波动方程，初始条件 $u_t^\varepsilon=h\equiv 0$ 且
\begin{align*}
u^\varepsilon=g^\varepsilon\coloneqq
\begin{cases}
e^{-\varepsilon/(r-2)^2}(r-1)(3-r) & \text{若 } 1<r<3,\\
0 & \text{其他},
\end{cases}
\end{align*}
其中 $r=|x|$。证明尽管 $|g^\varepsilon|\leqslant 1$，我们有
\begin{align*}
\max_{\mathbb{R}^2\times[0,4]}|u^\varepsilon|\to\infty \quad \text{当 } \varepsilon\to 0.
\end{align*}
（提示：使用 \S2.4.1\eqref{align:ch02-poisson-wave} 计算 $t>3$ 时 $u^\varepsilon(0,t)$。）
\end{exercise}

\begin{exercise}[波动方程的 Kelvin 变换。]\label{exercise:ch12-Kelvin}函数 $u:\mathbb{R}^n\times\mathbb{R}\to\mathbb{R}$ 的双曲 Kelvin 变换 $\mathcal{K}u=\bar{u}$ 为
\begin{align*}
\bar{u}(x,t)\coloneqq u(\bar{x},\bar{t})\,|||\bar{x}|^2-\bar{t}^2||^{(n-1)/2}=u\left(\frac{x}{|x|^2-t^2},\frac{t}{|x|^2-t^2}\right)\frac{1}{||x|^2-t^2|^{(n-1)/2}},
\end{align*}
只要 $|x|^2\neq t^2$，其中
\begin{align*}
\bar{x}=\frac{x}{|x|^2-t^2},\quad\bar{t}=\frac{t}{|x|^2-t^2}.
\end{align*}
证明若 $\square u=0$，则 $\square\bar{u}=0$。（与\cref{exercise:ch02-11} 比较。）
\end{exercise}

\begin{exercise}\label{align:ch12-sineGordon}
假设 $u$ 和 $v$ 求解方程组
\begin{align*}
\begin{cases}
(u-v)_t=2a\sin\left(\frac{u+v}{2}\right),\\
(u+v)_x=2a\sin\left(\frac{u-v}{2}\right),
\end{cases}
\end{align*}
其中 $a\neq 0$。证明 $w\coloneqq u$ 和 $w\coloneqq v$ 都求解 sine-Gordon 方程形式
\begin{align*}
w_{xt}=\sin w.
\end{align*}
为什么这等价于 PDE $\square w=\sin w$？
\end{exercise}

\begin{exercise}[续上题。]\label{exercise:ch12-Backlund}给定 sine-Gordon 方程的解 $v$，可尝试求解方程组\eqref{align:ch12-sineGordon} 来构造第二个解 $u$。此过程称为 B\"acklund 变换。从平凡解 $v\equiv 0$ 开始，使用 B\"acklund 变换对参数 $a$ 的每个选择计算另一个解 $u$。（提示：首先证明 $u$ 必须形如 $f(at+x/a)$。也证明 $(\tan(u/4))_t=a\tan(u/4)$。）
\end{exercise}

\begin{exercise}\label{exercise:ch12-SobolevHk}
证明 \S12.2.2 中 Sobolev 型不等式\eqref{align:ch12-Hk乘积估计}。（提示：若 $u\in H_k$ 且 $\beta$ 是多重指标 $|\beta|\leqslant k$，有估计 $\|D^\beta u\|_{L^p}\leqslant C\|u\|_{H_k}$，其中 $p$ 满足 (a) $p=\infty$ 若 $\frac{1}{2}-\frac{k}{n}+\frac{|\beta|}{n}<0$；(b) $2\leqslant p<\infty$ 若 $\frac{1}{2}-\frac{k}{n}+\frac{|\beta|}{n}=0$；(c) $\frac{1}{p}=\frac{1}{2}-\frac{k-|\beta|}{n}$ 若 $\frac{1}{2}-\frac{k}{n}+\frac{|\beta|}{n}>0$。假设多重指标 $|\beta_1|,\cdots,|\beta_r|$ 满足 (a)，$|\beta_{r+1}|,\cdots,|\beta_s|$ 满足 (b)，$|\beta_{s+1}|,\cdots,|\beta_m|$ 满足 (c)。估计 $\|D^{\beta_1}u_1\cdots D^{\beta_m}u_m\|_{L^2}$。）
\end{exercise}

\begin{exercise}\label{exercise:ch12-波映射}
光滑函数 $\mathbf{u}:\mathbb{R}^n\times[0,\infty)\to\mathbb{R}^m$，$\mathbf{u}=(u_1,\cdots,u_m)$，称为到单位球面 $S^{m-1}=\partial B(0,1)\subset\mathbb{R}^m$ 的\textbf{波映射}，若在 $\mathbb{R}^n\times[0,\infty)$ 中处处 $|\mathbf{u}|=1$ 且 $\mathbf{u}_{tt}-\Delta\mathbf{u}$ 在 $\mathbf{u}$ 处垂直于 $S^{m-1}$。证明因此 $\mathbf{u}$ 求解 PDE 方程组
\begin{align*}
\mathbf{u}_{tt}-\Delta\mathbf{u}=(|D\mathbf{u}|^2-|\mathbf{u}_t|^2)\mathbf{u}.
\end{align*}
\end{exercise}

\begin{exercise}\label{exercise:ch12-波映射能量}
证明若 $\mathbf{u}$ 是到单位球面的波映射，具有空间紧支集，则有能量守恒：
\begin{align*}
\frac{\mathrm{d}}{\mathrm{d}t}\int_{\mathbb{R}^n}|\mathbf{u}_t|^2+|D\mathbf{u}|^2\,\mathrm{d}x=0.
\end{align*}
\end{exercise}

\begin{exercise}[$p=5$ 时的小能量。]\label{exercise:ch12-小能量p5}改编 \S12.3.3\cref{theorem:ch12-次临界存在} 的证明以表明存在光滑解
\begin{align*}
\begin{cases}
u_{tt}-\Delta u+u^5=0 & \text{在 } \mathbb{R}^3\times(0,\infty) \text{ 中},\\
u=g,\;u_t=h & \text{在 } \mathbb{R}^3\times\{t=0\} \text{ 上},
\end{cases}
\end{align*}
只要能量 $E(0)$ 足够小。（提示：通过引入 $g$ 的 $L^6$ 范数修改 \S12.3.3 估计 (24)。）
\end{exercise}

为半线性波动方程发展的许多技术在非线性 Schr\"odinger（NLS）方程中有对应，我们将其余习题 devoted 于此。

\begin{exercise}\label{exercise:ch12-NLSGalilean}
设 $u$ 是非线性 Schr\"odinger 方程
\begin{align}
iu_t+\Delta u=f(|u|^2)u
\label{align:ch12-NLS}
\end{align}
在 $\mathbb{R}^n\times(0,\infty)$ 中的复值解。这里 $f:\mathbb{R}\to\mathbb{R}$。证明若 $\xi\in\mathbb{R}^n$，则
\begin{align*}
w(x,t)\coloneqq e^{\frac{i}{4}(2\xi\cdot x-|\xi|^2t)}u(x-\xi t,t)
\end{align*}
也求解 NLS 方程。这表明解的 Galileo 不变性。
\end{exercise}

\begin{exercise}\label{exercise:ch12-NLS守恒}
假设 $u$ 求解\cref{exercise:ch12-NLSGalilean}中的非线性 Schr\"odinger 方程\eqref{align:ch12-NLS}，且当其导数一起在 $|x|\to\infty$ 时快速衰减。推导以下恒等式：
\begin{align*}
\frac{\mathrm{d}}{\mathrm{d}t}\int_{\mathbb{R}^n}|u|^2\,\mathrm{d}x=0 \quad (\text{质量守恒}),
\end{align*}
\begin{align*}
\frac{\mathrm{d}}{\mathrm{d}t}\int_{\mathbb{R}^n}|Du|^2+F(|u|^2)\,\mathrm{d}x=0 \quad (\text{能量守恒}),
\end{align*}
其中 $F'=f$，以及
\begin{align*}
\frac{\mathrm{d}}{\mathrm{d}t}\int_{\mathbb{R}^n}\frac{\bar{u}Du-uD\bar{u}}{2i}\,\mathrm{d}x=0 \quad (\text{动量守恒}).
\end{align*}
记住 $|u|^2=u\bar{u}$。
\end{exercise}

\begin{exercise}[续上题。]\label{exercise:ch12-NLS爆破}
\begin{enumerate}[(1)]
\item 在前一题假设下，推导恒等式
\begin{align*}
\frac{\mathrm{d}^2}{\mathrm{d}t^2}\int_{\mathbb{R}^n}|x|^2|u|^2\,\mathrm{d}x=8\int_{\mathbb{R}^n}|Du|^2\,\mathrm{d}x+4n\int_{\mathbb{R}^n}f(|u|^2)|u|^2-F(|u|^2)\,\mathrm{d}x.
\end{align*}
\item 使用 (1) 证明不存在三次 NLS 方程
\begin{align*}
iu_t+\Delta u+|u|^2u=0
\end{align*}
在 $\mathbb{R}^n\times(0,\infty)$ 中对所有 $t\geqslant 0$ 存在的解，若
\begin{align*}
E(0)=\int_{\mathbb{R}^n}|Du(x,0)|^2-\frac{|u(x,0)|^4}{2}\,\mathrm{d}x<0
\end{align*}
且 $n\geqslant 2$。（R. Glassey, \textit{J. Math. Physics} 18 (1977), 1794--1797）
\end{enumerate}
\end{exercise}

\end{document}